\chapter{Master Problem Bank: Thermodynamics of Cells \& Nernst Equation}
\label{chap:qb_thermodynamics}

\begin{tcolorbox}[enhanced,colback=subtleamber,colframe=accentamber,arc=2mm,boxrule=1pt,
    title=\textbf{\large Part III Organization: Nernst Equation \& Cell Thermodynamics}]
This part compiles all questions and multi-step quantitative problems on the Nernst equation, equilibrium constants, cell Gibbs energy, enthalpy, entropy, temperature coefficients, and reversible cell heat from all 8 books. Identical and redundant variants from other books are co-located beneath each primary problem with full source citations.
\end{tcolorbox}

\section{Subtopic 3.1: Nernst Equation Calculations}

\begin{problembox}
\textbf{Q.3.1} \hfill \textbf{[Primary Source: Neeraj Kumar Part 1 Ex-I Q42 / NA Level-1 Q105]}
\label{q:qb-c3-01}

Calculate the EMF of the cell at $25^\circ\text{C}$:
\begin{equation*}
    \ce{Zn(s)} \mid \ce{Zn^2+(aq, } 0.01\,\text{M}\ce{)} \parallel \ce{Fe^2+(aq, } 0.001\,\text{M}\ce{)} \mid \ce{Fe(s)}
\end{equation*}
Given: $E^\circ(\ce{Zn^2+/Zn}) = -0.763\,\text{V}$ and $E^\circ(\ce{Fe^2+/Fe}) = -0.440\,\text{V}$.
\begin{multicols}{2}
\begin{enumerate}[label=(\alph*)]
    \item $+0.3525\,\text{V}$
    \item $+0.2935\,\text{V}$
    \item $+0.3230\,\text{V}$
    \item $-0.3230\,\text{V}$
\end{enumerate}
\end{multicols}

\tcbline
\textbf{\color{accentamber}Cross-Book Redundant / Identical Variants:}
\begin{itemize}[leftmargin=*,itemsep=2pt]
    \item \textbf{[Variant v1: Pearson Ex-I Q45, p. 4]} For the cell $\ce{Cd} \mid \ce{Cd^2+}(0.01\,\text{M}) \parallel \ce{Ag+}(0.5\,\text{M}) \mid \ce{Ag}$, given $E^\circ(\ce{Cd^2+/Cd}) = -0.40\,\text{V}$ and $E^\circ(\ce{Ag+/Ag}) = +0.80\,\text{V}$, calculate $E_{\text{cell}}$ at $298\,\text{K}$.
    \item \textbf{[Variant v2: GRB Practice Problem 12, p. 47]} Calculate the potential of a Daniell cell when $[\ce{Zn^2+}] = 0.5\,\text{M}$ and $[\ce{Cu^2+}] = 0.01\,\text{M}$. Standard potential $E^\circ_{\text{cell}} = 1.10\,\text{V}$.
    \item \textbf{[Variant v3: Cengage DPP 3.3 Q3, p. 5]} The potential of a hydrogen electrode at $298\,\text{K}$ in contact with a solution of $\text{pH} = 3$ and $P_{\ce{H2}} = 1\,\text{atm}$ is: (a) $-0.177\,\text{V}$ (b) $+0.177\,\text{V}$ (c) $-0.059\,\text{V}$ (d) $0.000\,\text{V}$.
\end{itemize}
\end{problembox}

\begin{problembox}
\textbf{Q.3.2} \hfill \textbf{[Primary Source: Narendra Avasthi Level-2 Q38 / PRSN Ex-II Q18]}
\label{q:qb-c3-02}

For the cell $\ce{Pt} \mid \ce{H2}(1\,\text{bar}) \mid \ce{HA}(0.1\,\text{M}) \parallel \ce{Ag+}(0.1\,\text{M}) \mid \ce{Ag}$, the cell EMF is measured to be $0.985\,\text{V}$ at $298\,\text{K}$. Given $E^\circ(\ce{Ag+/Ag}) = +0.800\,\text{V}$, calculate the acid dissociation constant $K_a$ of the weak monobasic acid \ce{HA}.

\tcbline
\textbf{\color{accentamber}Cross-Book Redundant / Identical Variants:}
\begin{itemize}[leftmargin=*,itemsep=2pt]
    \item \textbf{[Variant v1: Neeraj Kumar Part 1 Ex-II Sec A Q24, p. 14]} The EMF of the cell $\ce{Pt, H2}(1\,\text{atm}) \mid \text{Buffer } (\ce{CH3COOH} + \ce{CH3COONa}) \parallel \ce{KCl}(\text{sat}) \mid \ce{Hg2Cl2}, \ce{Hg}$ is $0.518\,\text{V}$. If $E_{\text{SCE}} = 0.242\,\text{V}$, calculate the pH of the buffer.
    \item \textbf{[Variant v2: GRB Practice Problem 24, p. 48]} A hydrogen electrode dipping into a solution of an organic acid of unknown strength gives an EMF of $0.236\,\text{V}$ when coupled with a standard calomel electrode. Calculate the pH and $[\ce{H+}]$ of the solution.
\end{itemize}
\end{problembox}

\section{Subtopic 3.2: Equilibrium Constant ($K_{eq}$) from $E^\circ_{\text{cell}}$}

\begin{problembox}
\textbf{Q.3.3} \hfill \textbf{[Primary Source: GRB Practice Problem 34, p. 49 / NA Level-1 Q118]}
\label{q:qb-c3-03}

Find the equilibrium constant $K_{eq}$ at $298\,\text{K}$ for the reaction:
\begin{equation*}
    \ce{Cu^2+(aq) + In+(aq) <=> Cu+(aq) + In^2+(aq)}
\end{equation*}
Given: $E^\circ(\ce{Cu^2+/Cu+}) = +0.15\,\text{V}$ and $E^\circ(\ce{In^2+/In+}) = -0.42\,\text{V}$.
\begin{multicols}{2}
\begin{enumerate}[label=(\alph*)]
    \item $1.0 \times 10^5$
    \item $4.4 \times 10^9$
    \item $2.1 \times 10^{19}$
    \item $8.8 \times 10^{-10}$
\end{enumerate}
\end{multicols}

\tcbline
\textbf{\color{accentamber}Cross-Book Redundant / Identical Variants:}
\begin{itemize}[leftmargin=*,itemsep=2pt]
    \item \textbf{[Variant v1: Neeraj Kumar Part 1 Ex-I Q65, p. 6]} Calculate the equilibrium constant for the disproportionation reaction $\ce{2Cu+(aq) <=> Cu^2+(aq) + Cu(s)}$ at $25^\circ\text{C}$. Given $E^\circ(\ce{Cu+/Cu}) = +0.52\,\text{V}$ and $E^\circ(\ce{Cu^2+/Cu}) = +0.34\,\text{V}$.
    \item \textbf{[Variant v2: Pearson Ex-I Q60, p. 6]} For the cell reaction $\ce{Sn(s) + 2Fe^3+(aq) -> Sn^2+(aq) + 2Fe^2+(aq)}$, $E^\circ_{\text{cell}} = 0.61\,\text{V}$. The equilibrium constant $K$ is approximately: (a) $10^{10}$ (b) $10^{20}$ (c) $10^{30}$ (d) $10^5$.
    \item \textbf{[Variant v3: Cengage DPP 3.3 Q7, p. 5]} The relationship between $\Delta G^\circ$ and $E^\circ_{\text{cell}}$ implies that for an electrochemical cell to be in chemical equilibrium, the required condition is: (a) $E^\circ_{\text{cell}} = 0$ (b) $E_{\text{cell}} = 0$ (c) $\Delta G^\circ = 0$ (d) $Q = 1$.
\end{itemize}
\end{problembox}

\section{Subtopic 3.3: Cell Thermodynamics: $\Delta G, \Delta S, \Delta H$ \& Temperature Coefficient}

\begin{problembox}
\textbf{Q.3.4} \hfill \textbf{[Primary Source: Atkins Topic 6C.4 / Neeraj Kumar Part 1 Ex-II Sec A Q32]}
\label{q:qb-c3-04}

The EMF of the standard Weston cadmium cell is expressed as a function of temperature:
\begin{equation*}
    E_t = 1.01845 - 4.05 \times 10^{-5}(t - 20) - 9.5 \times 10^{-7}(t - 20)^2 \quad (\text{in Volts})
\end{equation*}
where $t$ is the temperature in $^\circ\text{C}$. For the cell reaction:
\begin{equation*}
    \ce{Cd(s) + Hg2SO4(s) + \tfrac{8}{3}H2O(l) -> CdSO4.\tfrac{8}{3}H2O(s) + 2Hg(l)} \quad (n = 2)
\end{equation*}
Calculate at $t = 25^\circ\text{C}$:
\begin{enumerate}[label=(\alph*)]
    \item The temperature coefficient $\left(\frac{\partial E}{\partial T}\right)_P$
    \item The Gibbs free energy change $\Delta G$ in $\text{kJ\,mol}^{-1}$
    \item The entropy change $\Delta S$ in $\text{J\,K}^{-1}\,\text{mol}^{-1}$
    \item The reaction enthalpy $\Delta H$ in $\text{kJ\,mol}^{-1}$
    \item The reversible heat exchange $q_{\text{rev}}$ with the surroundings
\end{enumerate}

\tcbline
\textbf{\color{accentamber}Cross-Book Redundant / Identical Variants:}
\begin{itemize}[leftmargin=*,itemsep=2pt]
    \item \textbf{[Variant v1: GRB Practice Problem 28, p. 48]} The EMF of a galvanic cell is $1.05\,\text{V}$ at $298\,\text{K}$ and its temperature coefficient is $-4.0 \times 10^{-4}\,\text{V\,K}^{-1}$. If $n = 2$, calculate $\Delta H$ for the cell reaction.
    \item \textbf{[Variant v2: Narendra Avasthi Level-2 Q48, p. 28]} A cell has $E = 1.20\,\text{V}$ and $\Delta H = -200\,\text{kJ\,mol}^{-1}$ for a $2$-electron process. The temperature coefficient $(\partial E/\partial T)_P$ is: (a) $+5.4 \times 10^{-4}\,\text{V\,K}^{-1}$ (b) $-5.4 \times 10^{-4}\,\text{V\,K}^{-1}$ (c) $+2.7 \times 10^{-4}\,\text{V\,K}^{-1}$ (d) zero.
    \item \textbf{[Variant v3: Cengage DPP 3.4 Q4, p. 7]} If a cell absorbs heat from its surroundings during reversible operation at constant $T$ and $P$, then: (a) $(\partial E/\partial T)_P > 0$ (b) $(\partial E/\partial T)_P < 0$ (c) $\Delta H > \Delta G$ (d) Both (a) and (c).
\end{itemize}
\end{problembox}
